\section{A global flow and scalar comparison by localization}
\label{sec:flow}

We first construct the flow with estimates on fixed initial compact
sets.  Completeness of its positive-time slices will be proved in
Proposition~\ref{prop:harnack-completeness}.

\begin{theorem}[Localized flow]\label{thm:localized-flow}
There is a smooth K\"ahler--Ricci flow $h_t$ on $M$, for
$0\le t<\infty$, with $h_0=g$.  There are smooth real functions $U,P$
and $\rho$ such that
\begin{equation}\label{flow:identities}
 \begin{gathered}
 0<h_t\le g,\qquad
 h_t=g-t\Ric(g)+\dd U,\qquad U(\cdot,0)=0,\\
 P=U_t=\log\frac{\det h_t}{\det g},\qquad
 \rho=-P\ge0,\qquad R_{h_t}=-P_t\ge0.
 \end{gathered}
\end{equation}
For each finite $T$, there is $C_T$ such that
\begin{equation}\label{flow:rho-average}
 \sup_{0\le t\le T}
 \frac{1}{\Vol_gB_g(o,L)}\int_{B_g(o,L)}\rho(x,t)\,dV_g(x)
 \le C_TL\qquad(L\ge1).
\end{equation}
At every fixed $x$, $\rho(x,t)\le C_x(1+t)$ for all $t\ge0$.
Furthermore, on each finite time interval there is a smooth positive
function $Q$ and a constant $c>0$ with
\begin{equation}\label{flow:proper-supercaloric}
 Q\ge c(1+r^2),\qquad
 (\partial_t-\Delta_{h_t})Q\ge c\tr_{h_t}g.
\end{equation}
Smoothness at $t=0$ means smoothness with all one-sided time derivatives.
\end{theorem}

\subsection{Complete Hermitian approximations}

Take regular values $L\to\infty$ of the exhaustion $u$ in
Proposition~\ref{prop:initial-data}, and put $\Omega_L=\{u<L\}$.
Choose a fixed smooth nonnegative function $F$ on $[0,1)$, zero on
$[0,1/2]$, which equals $-2\log(1-v)$ near $v=1$.  On $\Omega_L$ set
\[
 f=F(u/L),\qquad \alpha=e^fg.
\]
The estimates for $u$ imply
\begin{equation}\label{flow:conformal-bounds}
 |\partial f|_\alpha\le C/L,
 \qquad |\dd f|_\alpha\le C/L^2.
\end{equation}
For example, near the boundary $e^f=(1-u/L)^{-2}$, so the factors
$(1-u/L)^{-1}$ and $(1-u/L)^{-2}$ in the first and second derivatives
are canceled by the conformal norm.  On the transition region
$u\asymp L$, the bound $|\dd u|_g\le C/u$ gives the same estimates.

The metric $\alpha$ is complete on $\Omega_L$: approaching the regular
boundary has infinite length for the factor $(1-u/L)^{-1}$, and the
closure of $\Omega_L$ is compact.  For each fixed $L$ it has bounded
geometry of every order in holomorphic charts.  Indeed choose finitely
many ordinary charts on the compact closure and, near a point of
boundary distance comparable to $1-u/L$, rescale their coordinate
radii by a sufficiently small fixed multiple of $1-u/L$.  In the
rescaled charts $\alpha$ and its inverse are uniformly bounded, and
every derivative of its coefficients is bounded.  The constants here
are allowed to depend on $L$ and on each derivative order.

Run Chern--Ricci flow on this complete Hermitian background:
\begin{equation}\label{flow:approximation}
 K=\alpha-t\Ric(\alpha)+\dd U,\qquad
 U_t=P=\log\frac{\det K}{\det\alpha},\qquad U(\cdot,0)=0.
\end{equation}
Until the passage to the local limit below, $U,P$ refer to the
approximating solution.
Bounded-geometry short-time existence is
\cite[Definition~2.1 and Lemma~2.1]{LeeTam}.  We establish estimates
on every closed bounded-geometry interval of this solution; the
resulting continuation intervals increase to infinity as $L$ increases.

Write $\mathscr L_K=\partial_t-\Delta_K$.
Differentiation of \eqref{flow:approximation} gives
\begin{equation}\label{flow:scalar-evolutions}
 \mathscr L_KP=-\tr_K\Ric(\alpha),\qquad
 \mathscr L_KP_t=-|\Ric(K)|_K^2,
 \qquad P_t(\cdot,0)\le C/L^2.
\end{equation}
The last inequality follows from
$\Ric(\alpha)=\Ric(g)-n\dd f$, initial Ricci positivity, and
\eqref{flow:conformal-bounds}.  In particular,
\begin{equation}\label{flow:upper-potentials}
 P_t\le C/L^2,
 \qquad P\le Ct/L^2,
 \qquad U\le Ct^2/(2L^2).
\end{equation}
These are bounded maximum-principle estimates on an already existing
bounded-geometry interval.

Fix a horizon $D>0$.  Suppose that $\phi\le0$ is smooth and
\begin{equation}\label{flow:phi-condition}
 \alpha-D\Ric(\alpha)+D\dd\phi\ge c\alpha
 \quad(c>0).
\end{equation}
Define
\begin{equation}\label{flow:Q-definition}
 Q=(D-t)P+U-D\phi+nt.
\end{equation}
Using
$\mathscr L_KU=P-n+\tr_K(\alpha-t\Ric(\alpha))$ gives the exact
identity
\begin{equation}\label{flow:Q-comparison}
 \mathscr L_KQ
 =\tr_K\{\alpha-D\Ric(\alpha)+D\dd\phi\}
 \ge c\tr_K\alpha,
 \qquad Q\ge0.
\end{equation}
The initial value is $-D\phi\ge0$.  For bounded $\phi$, the maximum
principle proves the last assertion directly; adding a vanishing
positive multiple of the proper function $f$ is an equivalent
localization.  The same proof applies to functions with downward
logarithmic poles, since $Q\to+\infty$ at those poles.

For a fixed $L$, choose finitely many canonical sections with no common
zero on $\overline{\Omega_L}$.  A constant shift of
$\log\sum_i|S_i|_g^2$ is bounded, nonpositive, and has complex Hessian
at least $\Ric(g)$.  It satisfies
\[
 \alpha-D\Ric(\alpha)+D\dd\phi
 \ge \alpha+nD\dd f\ge(1-CD/L^2)\alpha.
\]
Thus for fixed $D$ and sufficiently large $L$ we may take $c=1/2$.
Equations~\eqref{flow:upper-potentials} and
\eqref{flow:Q-comparison} give both bounds for $P$ on $0\le t\le D/2$:
indeed $(D-t)P\ge D\phi-U-nt$.  Integrating also gives both bounds
for $U$.  These bounds may depend on $L$.

\subsection{The largest metric eigenvalue}

We next bound the largest eigenvalue of $K$ relative to $\alpha$.
The bound will tend to one as $L\to\infty$ wherever $Q$ stays bounded
independently of $L$.  The local potential estimates below will supply
this control on each fixed compact set, yielding $h_t\le g$ in the
limit.  At a point use $g$-normal holomorphic
coordinates diagonalizing $K$, write $K_{i\bar i}=\lambda_i$, and
suppose that $\lambda_1$ is largest.  Put $q=K_{1\bar1}$ and
\[
 W=\frac{K_{1\bar1}}{e^fg_{1\bar1}}.
\]
The Rayleigh quotient in this fixed coordinate direction is a smooth
lower support for the largest relative eigenvalue and agrees with it
at the point.  It therefore suffices for a maximum-principle test of
that eigenvalue.  Wherever $W\ge1$, direct differentiation gives
\begin{equation}\label{flow:eigenvalue-raw}
 \mathscr L_K\log W
 \le -\frac1q\sum_{i,j}
            \frac{|\partial_1K_{i\bar j}|^2}{\lambda_i\lambda_j}
       +|\partial\log q|_K^2
       +CL^{-2}\tr_K\alpha.
\end{equation}
Here are the derivative exchanges in this formula.  Since $K-\alpha$
is closed,
\[
 \partial_1\partial_{\bar1}K_{i\bar i}
 -\partial_i\partial_{\bar i}K_{1\bar1}
 =\partial_1\partial_{\bar1}\alpha_{i\bar i}
 -\partial_i\partial_{\bar i}\alpha_{1\bar1}.
\]
In the chosen coordinates the right side is
$(e^f)_{1\bar1}-(e^f)_{i\bar i}$; the second derivatives of $g$
cancel by its K\"ahler symmetries.  Dividing their contraction by
$q\ge e^f$ bounds the error by $CL^{-2}\tr_K\alpha$.
Differentiating the denominator $g_{1\bar1}$ contributes
$-\sum_i\lambda_i^{-1}R_g(i,\bar i,1,\bar1)\le0$.
The remaining conformal derivatives have the same controlled size.
These facts give \eqref{flow:eigenvalue-raw}.

For $i\ne1$, the first-derivative exchange is
\[
 \partial_1K_{i\bar1}=\partial_iq-e^ff_i.
\]
The terms with $j=1$ in the negative square in
\eqref{flow:eigenvalue-raw} therefore dominate
\[
 \left|\partial\log q-W^{-1}\sum_{i\ne1}f_i\,dz_i\right|_K^2.
\]
Expanding this square and using
$\partial\log q=\partial\log W+\partial f$ at the point yields
\begin{equation}\label{flow:eigenvalue-gradient}
 \mathscr L_K\log W
 \le CL^{-2}\tr_K\alpha
       +2W^{-1}|\partial f|_K|\partial\log W|_K.
\end{equation}

Set $a=L^{-1}$ and
\[
 J(v)=av+\log(1+av)\qquad(v\ge0).
\]
At a positive maximum of $\log W-J(Q)$, one has
\[
 W\ge e^{aQ}(1+aQ),\qquad
 \partial\log W=J'(Q)\partial Q,
 \quad J'\ge a,\quad J'\le2a,\quad
 -J''=\frac{a^2}{(1+aQ)^2}.
\]
The chain rule, \eqref{flow:Q-comparison}, and
\eqref{flow:eigenvalue-gradient} give
\[
 \mathscr L_K(\log W-J(Q))
 \le (CL^{-2}-ca)\tr_K\alpha
 +2W^{-1}J'|\partial f|_K|\partial Q|_K
 +J''|\partial Q|_K^2.
\]
Young's inequality absorbs the middle term into half of
$-J''|\partial Q|_K^2$: its remaining coefficient is bounded by
\[
 C\frac{W^{-2}(J')^2}{-J''}|\partial f|_K^2
 \le C|\partial f|_K^2
 \le CL^{-2}\tr_K\alpha.
\]
For sufficiently large $L$, this contradicts the maximum principle.
To attain the maximum on $\Omega_L$, test first with the additional
penalty $-\delta f$.  Its Hessian term is controlled by
$C\delta L^{-2}\tr_K\alpha$, and its gradient terms have the same
bound after Young's inequality.  Send $\delta\downarrow0$ afterwards.
At $t=0$, $W=1$ and $Q\ge0$.  We conclude
\begin{equation}\label{flow:eigenvalue-bound}
 K\le e^{J(Q)}\alpha.
\end{equation}
With the bounded choice of $\phi$, $Q$ is bounded above by
\eqref{flow:upper-potentials}.  Thus \eqref{flow:eigenvalue-bound}
gives an upper metric bound.  The lower bound for $P$, hence for
$\det K/\det\alpha$, then bounds the least eigenvalue below.

We spell out the continuation implication in the uniform holomorphic
charts.  Let $\iota$ be the real symmetric representation of a Hermitian
matrix and write
\[
 F(B)=\tfrac12\log\det_{\R}B,\qquad
 S(x,t)=\iota(\alpha-t\Ric(\alpha)),\qquad
 T(D^2U)=\iota(\dd U).
\]
Then $U_t-F(S+T(D^2U))=-\log\det_{\C}\alpha$.
Two-sided metric equivalence puts the transformed Hessian $S+T(D^2U)$
in a fixed compact convex positive matrix range.  The function $F$
is smooth, concave and uniformly elliptic near this range.  The map
$T$ is linear, preserves nonnegativity, and has norm comparable to
the original norm on nonnegative real symmetric matrices.  The smooth
background coefficients have uniform local H\"older bounds.
These are the hypotheses of the transformed-Hessian parabolic
estimate \cite[Theorem~5.1]{ChuRegularity}.  Its bound for the spatial
second derivatives and first time derivative depends on $\|U\|_\infty$
and these data, without a prior bound for the full real Hessian.
Schauder estimates then give higher derivatives and continuation on
each fixed approximating background, as in
\cite[proof of Theorem~4.1]{LeeTam}.

For initial estimates, extend $U$ by zero to negative times and use
$S(x,t_+)$, where $t_+=\max(t,0)$.  Each fixed approximating solution
is already smooth one-sided at zero by short-time existence.
Since $U(\cdot,0)=U_t(\cdot,0)=0$, its spatial second derivatives
and first time derivative extend continuously by zero.  Thus the
extension is $C^2$ in the parabolic sense required by Theorem~5.1
of \cite{ChuRegularity} and solves the extended equation also at zero.
The time-Lipschitz background satisfies that theorem's H\"older
hypothesis, so the same estimate applies across the initial surface.
Spatial differentiation and the linear equation for $P$, with smooth
one-sided coefficients and initial data, give successive initial
Schauder estimates.  On each fixed compact these estimates are
uniform in $L$ once the local metric bounds below hold and $\alpha=g$
there.  Thus the approximating flows exist on $[0,D/2]$ whenever
$L$ is sufficiently large for the chosen $D$.

\subsection{The local limit and its potential estimates}

The bounds needed to pass to a limit use other choices of $\phi$.
For any canonical section $S$ from
Proposition~\ref{prop:initial-data}, take
\begin{equation}\label{flow:singular-weight}
 \phi=\log|S|_g^2-nf-\epsilon u^2-A_\epsilon,
 \qquad 0<\epsilon\le c_0/D,
 \qquad A_\epsilon=C_S'(1+\epsilon^{-1}).
\end{equation}
The growth bound for $S$ makes $\phi\le0$ after choosing $C_S'$.
Since $|\dd u^2|_g\le C$, the canonical-current inequality gives
\[
 \alpha-D\Ric(\alpha)+D\dd\phi
 \ge\alpha-CD\epsilon g\ge\tfrac12\alpha
\]
when $c_0$ is small.  The comparison is made away from the zeros of
$S$, where $\phi$ is smooth; the downward poles localize it away
from those zeros.  Finite sums of squared norms may be used as well.

Fix a compact set on which the chosen section does not vanish.
For all sufficiently large $L$ it lies in $\{u<L/2\}$, where $f=0$.
Equation~\eqref{flow:upper-potentials} bounds $Q$ above there by a
constant independent of $L$.  Equation~\eqref{flow:Q-comparison}
bounds $P$ below there, and \eqref{flow:eigenvalue-bound} gives a local
upper metric bound with $e^{J(Q)}\to1$.  The determinant lower bound
then gives a positive local lower metric bound.  Finitely many such
nonvanishing sections cover an arbitrary compact set.  Local
parabolic estimates on slightly larger compact sets give all
derivatives, including the one-sided initial estimates, uniformly
in $L$.  A diagonal subsequence over compacts and increasing horizons
converges smoothly to $U,P,h$ with $h_0=g$.  The approximations have
$\alpha=g$ on each fixed compact for large $L$, so
\eqref{flow:approximation} passes to the K\"ahler equation there.
The limit of \eqref{flow:eigenvalue-bound} is $h_t\le g$;
the limits of \eqref{flow:upper-potentials} give $P\le0$ and $P_t\le0$.
This proves \eqref{flow:identities}.

The $Q$ comparison also passes to the limit.  For a weight
$\phi=\log|S|_g^2-\epsilon u^2-A_\epsilon$ and $t\le D/2$, it gives,
using $U\le0$,
\begin{equation}\label{flow:rho-pointwise}
 \rho(x,t)\le
 \frac{-D\phi(x)+nt}{D-t}
 \le2\bigl(-\log|S(x)|_g^2+\epsilon u(x)^2+A_\epsilon+n\bigr).
\end{equation}
Take $D=2(T+1)$ and, on a ball of large radius $L$, take
$\epsilon$ comparable to $L^{-1}$ and at most $c_0/D$.
Proposition~\ref{prop:initial-data} then bounds the average of the
right side by $C_TL$.  Bounded radii are absorbed into the constant.
This proves \eqref{flow:rho-average}.  At a fixed $x$, choose $S(x)\ne0$,
put $D=2(t+1)$, and take $\epsilon=c_0/D$.  The same inequality gives
$\rho(x,t)\le C_x(1+t)$.

Finally replace $\log|S|_g^2$ by $\Xi$.  For every sufficiently small
$\epsilon>0$, the smooth function
$\phi_\epsilon=\Xi-\epsilon u^2-A_\epsilon$ is nonpositive after
a finite constant shift and gives a smooth nonnegative $Q_\epsilon$
with the differential inequality in \eqref{flow:Q-comparison}, now
with $\alpha=g$.  The difference
\[
 Q_\epsilon-Q_{\epsilon/2}
 =\frac{D\epsilon}{2}u^2+D(A_\epsilon-A_{\epsilon/2})
\]
shows that $Q_\epsilon\ge c u^2-C$ because
$Q_{\epsilon/2}\ge0$.  Adding a constant gives a positive $Q$
satisfying \eqref{flow:proper-supercaloric}.  Theorem~\ref{thm:localized-flow}
is proved.

\subsection{A cutoff adapted to the evolving metric}

The comparison arguments below require a positive integrable weight
$\chi$ with $\dd\chi\le C\chi h_t$.  A polynomial weight $u^{-k}$
gives only a Hessian bound by $Cu^{-k-2}g$, and $h_t$ need not be
uniformly comparable to $g$ at infinity.  We therefore construct a
function $H$ whose complex Hessian approaches $g-h_t$ at infinity.
The factor $\exp(-AH/u^2)$, for a suitable $A>0$, will contribute
the term $-Au^{-2}(g-h_t)$ to the logarithmic Hessian, absorbing
the errors measured in the initial metric.

Fix $T<\infty$ and work, if necessary, on a slightly larger finite
time interval.  Put
\[
 \eta_t=g-h_t,
 \qquad a_t=\tr_g\eta_t.
\]
The forms $\eta_t$ are nonnegative, closed, and bounded by $g$;
their traces satisfy $0\le a_t\le n$.
The potential identity gives
\[
 a_t=tR_g-\Delta_gU.
\]
Integrate against a cutoff from Section~\ref{sec:preliminaries} equal
to one on $B(o,L)$, and use
$-U=\int_0^t\rho(\cdot,s)\,ds$,
\eqref{prelim:averages}, and \eqref{flow:rho-average}.  Moving the
Laplacian onto the cutoff gives
\begin{equation}\label{flow:eta-average}
 \frac{1}{\Vol B(o,L)}\int_{B(o,L)}a_t\,dV_g
 \le C_Tt/L\qquad(L\ge1,\ 0\le t\le T).
\end{equation}

For large dyadic $L$, define
\[
 H_{L,t}=-\int_0^{L^2}H_ba_t\,db.
\]
There are $\delta,\beta>0$ such that, on $B(o,C_0L)$,
\begin{align}
 |H_{L,t}|+L|\partial H_{L,t}|_g
     &\le C_TL^{2-\delta}t^\beta,
                  \label{flow:heat-correction-size}\\
 \dd H_{L,t}&=\eta_t+O(t/L)g,
                  \label{flow:heat-correction-hessian}\\
 \partial_tH_{L,t}&\ge H_{L,t}/t-C_TL
                  \qquad(t>0).
                  \label{flow:heat-correction-time}
\end{align}
We give the heat estimates and the time identity separately.

Gaussian upper bounds and relative volume comparison, combined with
\eqref{flow:eta-average}, give, for some fixed $p>2$,
\begin{equation}\label{flow:heat-integrand}
 H_ba_t(x)\le
 C_T\min\left\{1,\frac{t}{L}\left(\frac{L}{\sqrt b}\right)^p\right\},
 \qquad x\in B(o,C_0L),\quad0<b\le L^2.
\end{equation}
For clarity, when $\sqrt b\ll L$ the ball $B(x,\sqrt b)$ may be
small relative to $B(o,L)$.  Relative volume comparison loses at
most a fixed power of $L/\sqrt b$; outside $B(o,CL)$, Gaussian
annuli absorb both that power and the polynomial volume growth.
The estimate $a_t\le n$ supplies the first term in the minimum.
For the gradient, use the heat-kernel estimate
\[
 |\nabla_xH_b(x,y)|
 \le \frac{C}{\sqrt b}\,
 \frac{\exp(-c\,d_g(x,y)^2/b)}
 {\sqrt{\Vol B(x,\sqrt b)\Vol B(y,\sqrt b)}}.
\]
It follows from \cite[Theorem~1.3, equation~(1.6)]{SoupletZhang}
and the Gaussian upper bound in \cite[Corollary~3.1]{LiYau},
absorbing the factor $1+d_g(x,y)^2/b$ into the Gaussian.
Integrating its absolute value against $a_t\ge0$ and using the same
relative-volume and annular estimates as above gives
$C_Tb^{-1/2}(t/L)(L/\sqrt b)^p$; increase the fixed $p>2$ if needed.
The separate bound $a_t\le n$ gives $Cb^{-1/2}$.
Thus the gradient obeys the minimum in
\eqref{flow:heat-integrand}, multiplied by $C/\sqrt b$.
Splitting both integrals at
$b_*=L^2(t/L)^{2/p}$ shows that their bounds are respectively
\[
 C_TL^{2-2/p}t^{2/p},\qquad
 C_TL^{1-1/p}t^{1/p}.
\]
After increasing the finite-time constant, we can take
$\delta=\beta=1/p$ in \eqref{flow:heat-correction-size}.
Lemma~\ref{lem:hodge-heat} gives
\[
 \dd H_{L,t}=\eta_t-\mathcal H_{L^2}\eta_t.
\]
The last form is nonnegative, and its trace is bounded by $C_Tt/L$
by \eqref{flow:heat-integrand}.  This proves
\eqref{flow:heat-correction-hessian}.

The trace identity also gives
\begin{equation}\label{flow:H-potential-identity}
 H_{L,t}=-t\int_0^{L^2}H_bR_g\,db+H_{L^2}U-U.
\end{equation}
To justify it, insert initial-metric cutoffs in
$\int_0^{L^2}H_b\Delta_gU\,db=H_{L^2}U-U$.
Both $U$ and $\Delta_gU=tR_g-a_t$ have polynomially bounded absolute
ball integrals, by \eqref{flow:rho-average},
\eqref{prelim:averages}, and $0\le a_t\le n$.
Gaussian bounds for the scalar heat kernel and its gradient, together
with the absolute Laplacian bounds for the cutoffs, make the shell
errors tend to zero.  This proves the identity distributionally and
then smoothly locally.  Differentiating it in $t$ is justified by
the same ball-integral bounds on $P=-\rho$; alternatively integrate
the identity over a time interval and then differentiate locally.
Subtracting $H_{L,t}/t$ gives
\[
 \partial_tH_{L,t}-H_{L,t}/t
 =H_{L^2}(P-U/t)-(P-U/t).
\]
Since $P_t\le0$, $P-U/t\le0$ and $0\le U/t-P\le\rho$.
Its heat average is at most $C_TL$, by
\eqref{flow:rho-average} and Gaussian annular summation.  The
unaveraged term has the favorable sign.  This proves
\eqref{flow:heat-correction-time}.

Choose a nonnegative smooth dyadic partition of unity $\zeta_L(u)$,
with scales $L\asymp u$, and put
$H=\sum_L\zeta_L(u)H_{L,t}$, using finitely many harmless initial
scales on a fixed compact.  The partition derivatives have sizes
$O(u^{-1})$ and $O(u^{-2})$ in gradient and complex Hessian.
Equations~\eqref{flow:heat-correction-size}--\eqref{flow:heat-correction-time}
give
\begin{equation}\label{flow:glued-correction}
 \begin{gathered}
 |H|\le C_Tu^{2-\delta}t^\beta,
 \qquad |\partial H|_g\le C_Tu^{1-\delta}t^\beta,\\
 \dd H=\eta_t+O(u^{-\delta}t^\beta+t/u)g,
 \qquad H_t\ge H/t-C_Tu.
 \end{gathered}
\end{equation}
These functions are smooth in space and continuously differentiable
in time, one-sided at $t=0$, which is all that the integration below
uses.  In fact \eqref{flow:H-potential-identity} gives their first
time derivatives in terms of $H_{L^2}P$ and local smooth functions.
The uniform absolute ball-integral bounds for $P$ give continuity of
that heat average, including at $t=0$, by local convergence and
uniform Gaussian tail estimates.  Spatial regularity follows from
the heat equation and \eqref{flow:heat-correction-hessian}.

\begin{lemma}[Integrable comparison cutoff]\label{flow:cutoff-lemma}
For every sufficiently large fixed integer $k$, there is a positive
function $\chi(x,t)$ on $M\times[0,T]$, smooth in space and continuously
differentiable in time one-sided at the initial surface, such that
\begin{equation}\label{flow:cutoff-bounds}
 \begin{gathered}
 C_T^{-1}u^{-k}\le\chi\le C_Tu^{-k},\qquad
 |\partial\chi|_g\le C_T\chi/u,\\
 \dd\chi\le C_T\chi h_t,
 \qquad \partial_t\chi\le C_T(1+t^{\beta-1})\chi
 \quad(t>0).
 \end{gathered}
\end{equation}
\end{lemma}

\begin{proof}
Take
\[
 \chi=u^{-k}\exp(-AH/u^2).
\]
For fixed $A$, \eqref{flow:glued-correction} gives the two-sided
size estimate and the gradient estimate.  Expanding
$\chi^{-1}\dd\chi=\dd\log\chi+
\partial\log\chi\otimes\bar\partial\log\chi$ gives at infinity
\[
 \chi^{-1}\dd\chi
 \le u^{-2}\{-A(g-h_t)+C_kg+o(1)g\},
\]
uniformly on the fixed time interval.  The terms containing
$H/u^2$ or its gradient are $o(u^{-2})g$ after $A$ is fixed.
Choose $A>C_k+1$ and then take $u$ sufficiently large.  The negative
multiple of $g$ absorbs the remaining error and leaves
$\chi^{-1}\dd\chi\le Au^{-2}h_t$.  On the remaining compact,
smoothness and positivity of $h_t$ give the asserted bound.  Finally,
$\chi_t/\chi=-AH_t/u^2$ and \eqref{flow:glued-correction} give
\[
 \chi_t/\chi\le A|H|/(tu^2)+C_T/u
                  \le C_T(1+t^{\beta-1}).
\]
The exponent $\beta>0$ makes this time coefficient integrable at zero.
\end{proof}

\subsection{Minimality and nonlinear scalar comparison}

\begin{proposition}[Scalar comparisons]\label{prop:scalar-comparison}
Fix a finite $T>0$.  The following conclusions hold for the flow of
Theorem~\ref{thm:localized-flow}.
\begin{enumerate}[label=\textup{(\roman*)}]
\item The function $\rho$ is the minimal nonnegative solution with
      zero initial values of
      \begin{equation}\label{flow:rho-source}
      (\partial_t-\Delta_{h_t})\rho
                        =\tr_{h_t}\Ric(g).
      \end{equation}
      There are nonnegative smooth supercaloric functions $\rho_j$ on
      $M\times[0,T]$ such that
      \begin{equation}\label{flow:rho-tails}
      0\le\rho_j\le\rho,\qquad
      \rho_j\ge\rho-C_j,\qquad
      \rho_j\longrightarrow0
      \text{ locally uniformly on }M\times[0,T].
      \end{equation}
      Each fixed $\rho_j$ has all spatial and one-sided time
      derivatives smooth at $t=0$ and bounded on compact spacetime
      sets.
\item Suppose $v$ is upper semicontinuous, bounded on compact
      spacetime sets, $0\le v\le\rho$, and satisfies, in the upper
      viscosity sense on $0<t<T$,
      \begin{equation}\label{flow:nonlinear-comparison}
      \partial_tv\le e^{-v}\Delta_{h_t}v+(1-e^{-v})R_{h_t}.
      \end{equation}
      Then $v=0$ on $M\times[0,T)$.
\end{enumerate}
Neither assertion requires completeness of $h_t$.
\end{proposition}

Before proving the proposition, we isolate the smooth majorant used
in part \textup{(ii)}.  Its construction uses compact-domain parabolic
theory and will be given immediately after the comparison proof.

\begin{lemma}[Smooth majorant]\label{lem:scalar-majorant}
Fix $T>0$ and let $v$ be upper semicontinuous, bounded on compact
spacetime sets, and satisfy $0\le v\le\rho$ on $M\times[0,T]$.
Suppose that $v$ satisfies \eqref{flow:nonlinear-comparison} in the
upper viscosity sense for $0<t<T$.  There is a function $w$, smooth
on $M\times(0,T)$ and continuous on $M\times[0,T)$, such that
\[
 v\le w\le\rho,\qquad
 w_t=e^{-w}\Delta_{h_t}w+(1-e^{-w})R_{h_t},\qquad
 w(\cdot,0)=0.
\]
\end{lemma}

\begin{proof}[Proof of Proposition~\ref{prop:scalar-comparison}]
The Ricci identity
$\Ric(h_t)=\Ric(g)+\dd\rho$, traced with $h_t$, proves
\eqref{flow:rho-source}, since $\rho_t=R_{h_t}$.
Its source is nonnegative.
Solve the source equation with zero initial and lateral boundary
data on a smooth relatively compact exhaustion, first as a weak
energy solution \cite[Chapter~III, Theorem~4.1]{LSU}, and then use
local parabolic regularity.  A source need not vanish at the
initial lateral corner, so full corner compatibility is not
asserted.  This bounded-domain construction and its energy
estimates pass between a finite set of coordinate charts; it
requires no completeness of the evolving metric.
Domain comparison makes these solutions increase, and
comparison with $\rho$ bounds them above.  Their limit $\rho_{\min}$
is the minimal nonnegative potential.  Local estimates give a smooth
solution for $t>0$ and continuous zero initial values, since it is
bounded by the smooth function $\rho=O_K(t)$ on each compact.
The difference $z=\rho-\rho_{\min}$ satisfies
\[
 0\le z\le\rho,\qquad z_t=\Delta_{h_t}z,
 \qquad z(\cdot,0)=0.
\]

Use $\chi$ from Lemma~\ref{flow:cutoff-lemma}, with $k>2n+3$,
and a cutoff $p_N=p(u/N)$ which is one on $u\le N$ and zero on
$u\ge2N$.  Set
\[
 I_N(t)=\int_M z\chi p_N\,dV_{h_t}.
\]
Since $\partial_tdV_{h_t}=-R_{h_t}dV_{h_t}$ and $R_{h_t}\ge0$,
integration by parts on the compact support of $p_N$ gives
\begin{equation}\label{flow:linear-mass}
 I_N'(t)\le C_T(1+t^{\beta-1})I_N(t)+E_N(t),
 \qquad |E_N(t)|\le C_TN^{2n-1-k}.
\end{equation}
Here the error contains the derivatives falling on $p_N$.  Its
complex Hessian, including the gradient cross terms with $\chi$,
has $g$-norm at most $C_TN^{-k-2}$ on $N\le u\le2N$.
The pointwise cofactor identity and $0<h_t\le g$ imply, for any
real $(1,1)$-form $B$,
\begin{equation}\label{flow:cofactor}
 |\tr_{h_t}B|\,dV_{h_t}
 \le C_n|B|_g\,dV_g.
\end{equation}
Indeed diagonalize $h_t$ relative to $g$; the coefficient multiplying
each diagonal entry of $B$ is the product of the other $n-1$
eigenvalues, each at most one.  This also covers $n=1$.
Equation~\eqref{flow:rho-average} bounds the integral of $z$ on the
shell by $C_TN^{2n+1}$, proving the error bound.

The mass $I=\int z\chi\,dV_h$ is finite.  Its tails tend to zero
uniformly in $t$ by $z\le\rho$, $dV_h\le dV_g$, and the same
dyadic shell estimate.  Its initial limit is zero by local
convergence and this tail control.  Integrating
\eqref{flow:linear-mass} from a positive initial time, sending
$N\to\infty$, and then sending that time to zero gives
$I(t)\le\int_0^tC_T(1+s^{\beta-1})I(s)\,ds$.
Gronwall's inequality gives $I=0$, and positivity of $\chi$ gives
$z=0$.  Thus $\rho=\rho_{\min}$.

For the tails, choose smooth spatial cutoffs $0\le\theta_j\le1$
increasing to one, each compactly supported and eventually equal
to one near each fixed compact.  Let $\rho_j$ be the minimal
zero-data potential of
\[
 F_j=(1-\theta_j)\tr_{h_t}\Ric(g).
\]
It is nonnegative and supercaloric.  Linearity on each exhaustion
domain and monotone limits give the decomposition of $\rho$ into
this exterior-source potential and the minimal potential of
$\theta_j\tr_h\Ric(g)$.  The latter is at most
$T\sup_{M\times[0,T]}\theta_j\tr_h\Ric(g)<\infty$ by the scalar
maximum principle on the domains.  This proves
$\rho_j\ge\rho-C_j$.  Monotone convergence of the nonnegative
Dirichlet Green integrals, first in the domains and then for the
source cutoffs, shows that $\rho_j\downarrow0$ pointwise.

We record the initial regularity needed later.  On every fixed
compact neighborhood the coefficients of $\Delta_{h_t}$ and the
source $F_j$ are smooth and uniformly parabolic through $t=0$.
The domination $0\le\rho_j\le\rho=O_K(t)$ supplies continuous
zero initial values.
For clarity, this initial regularity is local in space.  Extend
the smooth coefficients and source from a coordinate neighborhood
to a smooth uniformly parabolic Cauchy problem, and solve it with
zero initial data \cite[Chapter~IV, Theorem~5.1]{LSU}.
Subtract this local particular solution from $\rho_j$ on a smaller
neighborhood.  The difference is $O(t)$ there and solves the
homogeneous equation for $t>0$.  A spatially localized energy
estimate from time $\delta>0$, followed by $\delta\downarrow0$,
puts the difference in the local weak energy class through zero;
the squared $L^2$ norm of its initial value on the fixed compact
is $O(\delta^2)$ and tends to zero.  Extend that difference by zero to
negative times, and extend the coefficients smoothly.  Its zero
initial trace removes a distributional boundary term, so it
solves the homogeneous equation across the initial surface.
Interior regularity \cite[Chapter~III, Theorem~12.1]{LSU} makes it
smooth there.  Adding back the particular solution proves the
claimed smooth one-sided jets.  This argument extends the
homogeneous difference, not the forced potential, by zero.
In particular,
\[
 \partial_t\rho_j(x,0)=F_j(x,0),
\]
and higher initial derivatives are determined recursively from
$(\rho_j)_t=\Delta_{h_t}\rho_j+F_j$; they are smooth functions of
$x$.  No regularity at a distant Dirichlet corner is used.  Dini's
theorem on compact spacetime sets now gives the local uniform
convergence in \eqref{flow:rho-tails}.  If an auxiliary smooth
extension across $t=0$ is desired on a compact, it can be chosen
to match these jets.  Extension by zero to negative time is not
asserted.

For \textup{(ii)}, let $w$ be the smooth majorant from
Lemma~\ref{lem:scalar-majorant}.  Thus $v\le w\le\rho$, and $w$
solves equality in \eqref{flow:nonlinear-comparison}.
Multiplying its equation by $e^w$ gives the useful
volume-form identity
\begin{equation}\label{flow:nonlinear-volume}
 \partial_t\bigl[(e^w-1)dV_{h_t}\bigr]
                    =(\Delta_{h_t}w)dV_{h_t}.
\end{equation}
For example, the cancellation uses
$\partial_tdV_h=-R_hdV_h$ and
$e^w(1-e^{-w})R_h=(e^w-1)R_h$.

Integrate \eqref{flow:nonlinear-volume} against $\chi p_N$ and
move the Laplacian onto that compactly supported function.  Since
$w\le e^w-1$ and $\Delta_h\chi\le C_T\chi$, the interior term is
bounded by a constant times
\[
 M_N(t)=\int_M(e^w-1)\chi p_N\,dV_h.
\]
The time derivative of $\chi$ is bounded by
$C_T(1+t^{\beta-1})\chi$.  The exterior error is again
$O(N^{2n-1-k})$, using $w\le\rho$ and
\eqref{flow:cofactor}.  Moreover
\[
 0\le(e^w-1)dV_h\le e^\rho dV_h=dV_g.
\]
Thus the full weighted mass is finite, its tails are uniformly
small, and its initial limit is zero by dominated convergence on
compacts and $\chi\asymp u^{-k}$.  The same limiting Gronwall
argument as above forces it to vanish.  Hence $w=0$, and $0\le v\le w$
proves \textup{(ii)}.
\end{proof}

\subsection{Constructing the smooth majorant}

\begin{proof}[Proof of Lemma~\ref{lem:scalar-majorant}]
We construct the majorant on relatively compact domains and then
pass to a limit.
On a smooth relatively compact domain
$\Omega$ and for a fixed time less than $T$, consider equality in
\eqref{flow:nonlinear-comparison}.  Zero and $\rho$ are respectively
a subsolution and a supersolution.  To check the latter, use
\[
 R_{h_t}-\Delta_{h_t}\rho=\tr_{h_t}\Ric(g)\ge0.
\]
The function $\rho+2\alpha$ is also a supersolution for every
constant $\alpha>0$.

Take constant initial data $\alpha$.  On $\partial\Omega$ choose
a smooth $\zeta:[0,\infty)\to[0,1]$, equal to one near zero and
supported in $[0,1]$, and put
\[
 \varphi(x,t)=\rho(x,t)+\alpha
       -e^{-\alpha}R_g(x)t\zeta(t/\delta).
\]
For $\delta\sup_{\partial\Omega}R_g\le\alpha/2$ this lies between
$\rho+\alpha/2$ and $\rho+\alpha$.
It has $\varphi(\cdot,0)=\alpha$ and
$\varphi_t(\cdot,0)=(1-e^{-\alpha})R_g$, which are exactly the
value and first equation compatibility conditions.

We give the local existence and continuation argument on this
fixed compact domain.  We use the compatible linear Dirichlet
Schauder theorem \cite[Chapter~IV, Theorem~5.2]{LSU} and its
uniform small-time inverse estimate
\cite[Chapter~IV, Theorem~5.4]{LSU}.
These Euclidean-domain statements have the same compact-manifold
form by the following finite-chart construction for a smooth
uniformly parabolic linear operator $L$.  Choose a partition of
unity $\chi_i$ and cutoffs $\psi_i=1$ near $\supp\chi_i$ in
finitely many interior and boundary charts.  Local domains are
smoothly capped outside these supports, and $\psi_i$ vanishes near
every artificial boundary.  If $E_i$ is the local zero-data inverse,
then $Sf=\sum_i\psi_i E_i(\chi_i f)$ satisfies
$LSf=f+\mathcal Rf$ for the global linear operator $L$.
The remainder consists of first- and zeroth-order commutators
$[L,\psi_i]E_i(\chi_i f)$.  On the source space
$f(\cdot,0)=0$, the local solutions have zero initial trace.
Parabolic interpolation and the uniform inverse bounds give
$\|\mathcal R\|_{C^{\beta,\beta/2}\to C^{\beta,\beta/2}}
\le C\epsilon^\theta$ on a time interval of length $\epsilon$,
with $\theta>0$; for $0<\beta<1$ one can take
$\theta=(1-\beta)/2$.  The remainder preserves this source space.
Thus $S(I+\mathcal R)^{-1}$ is the required global inverse for
small $\epsilon$.  All constants here may depend on $\Omega$.

Extend the boundary data to a smooth function $p$ with
\[
 p(x,0)=\alpha,\qquad
 p_t(x,0)=c(x):=(1-e^{-\alpha})R_g(x)
 \quad\hbox{throughout }\Omega.
\]
For example, in a smooth boundary collar with projection $\pi$
and cutoff $\vartheta=1$ near the boundary, take
\[
 p(x,t)=\alpha+t c(x)
   +\vartheta(x)\{\varphi(\pi(x),t)-\alpha-tc(\pi(x))\},
\]
and use $\alpha+tc(x)$ outside the collar.  The bracket and its
first time derivative vanish at zero.
Set $L=\partial_t-e^{-p}\Delta_h$ and seek
$v_{\Omega,\alpha}=p+z$, with zero
initial and lateral data for $z$.  Its equation is
\[
 Lz=\mathcal N(z):=
  (e^{-p-z}-e^{-p})\Delta_h(p+z)
  +(1-e^{-p-z})R_h-Lp.
\]
For every such $z\in C^{2+\beta,1+\beta/2}$,
$\mathcal N(z)(x,0)=0$ throughout $\Omega$.
The linear inverse is therefore applied to compatible sources.
If $z(\cdot,0)=0$, then
$\|z\|_\infty\le\epsilon\|z_t\|_\infty$; spatial interpolation
and the time H\"older quotient give
\[
 \|z\|_{C^{\beta,\beta/2}}
 \le C\epsilon^{1-\beta/2}
       \|z\|_{C^{2+\beta,1+\beta/2}}.
\]
On a fixed ball in the latter space, the product and composition
estimates consequently give
\[
 \|\mathcal N(z)-\mathcal N(w)\|_{C^{\beta,\beta/2}}
 \le C\epsilon^{1-\beta/2}
       \|z-w\|_{C^{2+\beta,1+\beta/2}}.
\]
The coefficient of every highest-derivative difference has this
small factor.  Also $\mathcal N(0)$ is smooth and vanishes at
$t=0$, so its $C^{\beta,\beta/2}$ norm tends to zero with $\epsilon$.
The linear inverse and the contraction theorem give a local
Dirichlet solution.  Classical comparison with the barriers gives
$0\le v_{\Omega,\alpha}\le\rho+2\alpha$.

For continuation, write
$\Delta_h=a^{ij}\partial_i\partial_j+b^j\partial_j$ in fixed
real coordinates and set $q=e^{v_{\Omega,\alpha}}$.  Directly,
\[
 q_t-\partial_i\left(\frac{a^{ij}}q\,\partial_jq\right)
  =\frac{b^j-\partial_i a^{ij}}q\,\partial_jq+(q-1)R_h.
\]
The barriers bound $q$ above and away from zero.  The leading
matrix is uniformly elliptic, and the displayed drift and forcing
are bounded independently of $\nabla q$ on the fixed cylinder.
Boundary as well as interior H\"older estimates therefore apply
\cite[Chapter~III, Theorem~10.1]{LSU}; the exterior-measure
condition there holds for the smooth compact boundary.
Hence $v_{\Omega,\alpha}$ is H\"older up to the parabolic boundary.
Its original equation now has H\"older coefficients and source,
so \cite[Chapter~IV, Theorem~5.2]{LSU} gives a controlled
$C^{2+\gamma,1+\gamma/2}$ norm on the closed slab for some
$0<\gamma<1$.  In particular its second spatial derivatives are
bounded there.  Only this finite corner regularity is asserted
from the first compatibility condition.

On slabs separated from the initial surface, differentiated
boundary Schauder estimates give all higher bounds, since the
coefficients and lateral data are smooth.  At a hypothetical
positive maximal time the solution thus has smooth limiting
initial data satisfying the boundary equation jets.  The same
local construction, with those data and their first equation jet
in place of $\alpha$ and $c$, restarts it.  This proves existence
up to the prescribed finite horizon.  No constants near the
remote boundaries need be uniform as $\Omega$ exhausts $M$.

Viscosity comparison with this smooth solution gives
$v\le v_{\Omega,\alpha}$.  To see the only nonlinear point in that
comparison, at a positive interior maximum of
$e^{-Ct}(v-v_{\Omega,\alpha})$ the test has the same spatial Hessian
as $v_{\Omega,\alpha}$.  The difference of the right sides is bounded
by a constant times $v-v_{\Omega,\alpha}$, because the derivative
in the value variable is
$e^{-v}(R_h-\Delta_hv_{\Omega,\alpha})$, bounded on the compact slab.
Choose $C$ larger than this bound and use a strict increasing time
penalty.  The initial and lateral gaps are positive, so no boundary
maximum interferes.  This proves the comparison directly, consistently
with the usual viscosity framework \cite{CIL}.

Exhaust $M$, let $\alpha\downarrow0$, and use local uniformly
parabolic estimates to extract a smooth interior limit $w$ with
\[
 v\le w\le\rho,\qquad
 w_t=e^{-w}\Delta_hw+(1-e^{-w})R_h,
 \qquad w(\cdot,0)=0.
\]
The limit at the initial surface follows from the barriers on each
compact; no uniform estimates near the remote lateral boundaries
are required.
\end{proof}
